\chapter{Aplicaciones lineales}
\label{cap:aplicaciones}

Los espacios de este tema son sobre el mismo cuerpo $\K$.
Las fórmulas numéricas de dimensión se aplican en dimensión finita.

\section{Linealidad y determinación por una base}
\label{sec:linealidad}

Una aplicación $f:U\to V$ es \textbf{lineal}, o un homomorfismo vectorial,
si para $u,w\in U$ y $\alpha,\beta\in\K$,
\[
 f(\alpha u+\beta w)=\alpha f(u)+\beta f(w).
\]
Equivale a conservar la suma y el producto por escalares.
Se deducen $f(0)=0$, $f(-u)=-f(u)$ y la conservación de cualquier
combinación lineal finita. Una relación de dependencia entre vectores
produce la misma relación entre sus imágenes. La independencia,
en cambio, puede perderse.

Para comprobar linealidad se verifica la igualdad para vectores y escalares
arbitrarios. Para refutarla basta un contraejemplo. La condición $f(0)=0$
es necesaria, pero no suficiente.
Son lineales la identidad, la aplicación nula, las homotecias $u\mapsto\lambda u$,
la trasposición de matrices y las aplicaciones $X\mapsto AX$ con $A$ fija.
También lo son la derivación de polinomios y la proyección canónica
$\pi:V\to V/W$, $\pi(v)=v+W$.

Si $B=(u_1,\ldots,u_n)$ es una base de $U$, se pueden elegir arbitrariamente
$v_1,\ldots,v_n\in V$ y existe una única aplicación lineal con $f(u_i)=v_i$:
\[
 f(x_1u_1+\cdots+x_nu_n)=x_1v_1+\cdots+x_nv_n.
\]
Las imágenes de una base determinan completamente la aplicación.
Si se dan imágenes de vectores dependientes, estas deben respetar sus
relaciones; no se pueden elegir de forma independiente.

\subsection*{Operaciones y subespacios}

El conjunto $\Lin(U,V)$ de aplicaciones lineales es un espacio vectorial
con las operaciones punto a punto:
\[
 (f+g)(u)=f(u)+g(u),\qquad (\lambda f)(u)=\lambda f(u).
\]
Si $f:U\to V$ y $g:V\to W$ son lineales, $g\circ f$ también lo es:
primero actúa $f$ y después $g$.

La imagen de un subespacio $U'\subseteq U$ y la imagen recíproca de un
subespacio $V'\subseteq V$ son subespacios:
\[
 \begin{split}
 f(U')&=\{f(u):u\in U'\},\\
 f^{-1}(V')&=\{u\in U:f(u)\in V'\}.
 \end{split}
\]
Si $U'=L(u_1,\ldots,u_r)$, entonces
$f(U')=L(f(u_1),\ldots,f(u_r))$ y $\dim f(U')\le\dim U'$.
La notación $f^{-1}(V')$ para un conjunto no exige que exista una aplicación
inversa de $f$.

\section{Núcleo, imagen y tipos de aplicaciones}
\label{sec:nucleo-imagen}

El \textbf{núcleo} y la \textbf{imagen} de $f:U\to V$ son
\[
 \Ker(f)=\{u\in U:f(u)=0\},\qquad
 \Ima(f)=f(U)=\{f(u):u\in U\}.
\]
Son subespacios de $U$ y de $V$, respectivamente. Si $B=(u_1,\ldots,u_n)$
es una base de $U$, las imágenes $f(u_i)$ generan $\Ima(f)$.
La \textbf{fórmula de dimensiones} es
\[
 \boxed{\dim(U)=\dim\Ker(f)+\dim\Ima(f).}
\]
La idea de su demostración es extender una base del núcleo a una base
de $U$: las imágenes de los vectores añadidos forman una base de la imagen.
El \textbf{rango de $f$} es $\rg(f)=\dim\Ima(f)$.

\begin{center}
\begin{tabular}{@{}lll@{}}
\toprule
Tipo & Significado & Criterio \\
\midrule
Monomorfismo & Lineal e inyectiva & $\Ker(f)=\{0\}$. \\
Epimorfismo & Lineal y sobreyectiva & $\Ima(f)=V$. \\
Isomorfismo & Lineal y biyectiva & Ambas condiciones. \\
\bottomrule
\end{tabular}
\end{center}
La inyectividad equivale a conservar la independencia de las familias,
a que las imágenes de una base de $U$ sean base de $\Ima(f)$ y a
$\rg(f)=\dim(U)$. La sobreyectividad equivale a que las imágenes
de un sistema generador de $U$ generen $V$, o a $\rg(f)=\dim(V)$.
Si $U$ y $V$ tienen la misma dimensión finita, inyectividad,
sobreyectividad y biyectividad son equivalentes.

La inversa de un isomorfismo es lineal. Dos espacios vectoriales son
\textbf{isomorfos} si existe un isomorfismo entre ellos; en dimensión finita,
esto ocurre si y solo si tienen la misma dimensión sobre $\K$.
En particular, todo espacio de dimensión $n$ es isomorfo a $\K^n$ mediante
sus coordenadas en una base. El isomorfismo conserva la estructura
vectorial; no tiene por qué conservar otras operaciones de los objetos.

\subsection*{Primer teorema de isomorfía}
\label{sec:isomorfia}

Toda aplicación lineal induce el isomorfismo
\[
 \widetilde f:U/\Ker(f)\longrightarrow\Ima(f),\qquad
 u+\Ker(f)\longmapsto f(u).
\]
Está bien definido porque dos representantes difieren en un vector del
núcleo; es inyectivo por esa misma condición y sobreyectivo por definición
de imagen. En consecuencia, $U/\Ker(f)\simeq\Ima(f)$.

Si $\pi:U\to U/\Ker(f)$ es la proyección canónica y
$i:\Ima(f)\hookrightarrow V$ la inclusión, la
\textbf{descomposición canónica} es $f=i\circ\widetilde f\circ\pi$:
un epimorfismo, seguido de un isomorfismo y de un monomorfismo.

\section{Representación matricial}
\label{sec:matriz-lineal}

Sean $B=(u_1,\ldots,u_n)$ una base de $U$ y
$C=(v_1,\ldots,v_m)$ una base de $V$. La matriz de $f$ respecto de
estas bases es
\[
 \M_{BC}(f)=\bigl(\coord{f(u_1)}{C}\ \cdots\ \coord{f(u_n)}{C}\bigr)
 \in\M_{m\times n}(\K).
\]
Cada \emph{columna} contiene las coordenadas de la imagen de un vector
de la base de partida. La expresión analítica de $f$ es
\[
 \coord{f(u)}{C}=\M_{BC}(f)\coord{u}{B}.
\]
En $\K^n\to\K^m$, si no se indican otras bases se usan las canónicas.
Con las bases fijadas, la asignación $f\mapsto\M_{BC}(f)$ es un
isomorfismo vectorial entre $\Lin(U,V)$ y $\M_{m\times n}(\K)$;
por tanto, $\dim\Lin(U,V)=mn$.

\subsection*{Cálculo del núcleo, imagen y preimágenes}

Para $A=\M_{BC}(f)$:
\begin{itemize}
 \item \textbf{Núcleo:} se resuelve $AX=0$. Sus soluciones son coordenadas
       en $B$; una base del espacio de soluciones se traduce a vectores de $U$.
 \item \textbf{Imagen:} las columnas originales de $A$ generan sus coordenadas
       en $C$. Se extrae una base según la sección~\ref{sec:rango-vectores}.
 \item \textbf{Imagen de un subespacio:} se aplican $A$ a las coordenadas
       de sus generadores y se selecciona una base de las imágenes.
 \item \textbf{Imagen recíproca:} si $V'$ tiene ecuaciones $DY=0$ en $C$,
       se resuelve $DAX=0$ y se interpretan las soluciones en $B$.
\end{itemize}
Además,
\[
 \rg(f)=\rg(A),\qquad \dim\Ker(f)=n-\rg(A).
\]
Así, $f$ es inyectiva si $\rg(A)=n$, sobreyectiva si $\rg(A)=m$,
y un isomorfismo si $m=n$ y $A$ es invertible.

\subsection*{Operaciones y cambios de base}

Con las mismas bases para ambas aplicaciones,
\[
 \M_{BC}(\alpha f+\beta g)=\alpha\M_{BC}(f)+\beta\M_{BC}(g).
\]
Si $f:U\to V$, $g:V\to W$ y $B,C,D$ son las respectivas bases,
\[
 \M_{BD}(g\circ f)=\M_{CD}(g)\M_{BC}(f).
\]
La base intermedia debe coincidir en ambos factores.
La identidad de $V$ en bases distintas tiene como matriz
$\M_{BB'}(\mathrm{Id})=\M_{BB'}$, la matriz de cambio de base.

Si se cambian $B,C$ por $B',C'$, entonces
\[
 \boxed{\M_{B'C'}(f)=\M_{CC'}\M_{BC}(f)\M_{B'B}.}
\]
Se lee de derecha a izquierda: expresar la entrada en $B$, aplicar $f$,
y expresar la salida en $C'$. Los factores de cambio son invertibles;
por ello, el rango de una aplicación no depende de las bases elegidas.

\section{Endomorfismos y automorfismos}
\label{sec:endomorfismos}

Un \textbf{endomorfismo} es una aplicación lineal $f:V\to V$;
un \textbf{automorfismo} es un endomorfismo biyectivo.
Se escribe $\Lin(V)=\Lin(V,V)$. Para representar un endomorfismo
se usa una misma base $B$ en partida y llegada:
$\M_B(f)=\M_{BB}(f)$ es cuadrada.

Las matrices de un mismo endomorfismo en bases distintas son semejantes.
Si $A=\M_B(f)$, $A'=\M_{B'}(f)$ y $P=\M_{B'B}$, entonces
\[
 A'=P^{-1}AP.
\]
Recíprocamente, toda matriz semejante a $A$ representa el mismo
endomorfismo en alguna base. Se conecta así el cambio de base con la
semejanza de la sección~\ref{sec:inversas}.

Los automorfismos forman el \textbf{grupo general lineal} $\GL(V)$
para la composición. En una base se corresponde con el grupo
$\GL(n,\K)$ de matrices invertibles para el producto.
La identidad es el neutro y la matriz de $f^{-1}$ es $\M_B(f)^{-1}$.

\section{Proyecciones y simetrías}
\label{sec:proyecciones}

Si $V=U\oplus W$, cada vector se escribe de forma única $v=u+w$,
con $u\in U$ y $w\in W$. Se definen la \textbf{proyección} de base $U$
y dirección $W$, y la \textbf{simetría} de la misma base y dirección:
\[
 p(u+w)=u,\qquad s(u+w)=u-w.
\]
Ambas son endomorfismos. En una base formada primero por una base de $U$
y después por una de $W$, con $d=\dim(U)$,
\[
 \M_B(p)=\begin{pmatrix}I_d&0\\0&0\end{pmatrix},\qquad
 \M_B(s)=\begin{pmatrix}I_d&0\\0&-I_{n-d}\end{pmatrix}.
\]
Para obtener sus matrices en otra base se aplica la fórmula de semejanza.

\subsection*{Caracterizaciones}

Una proyección satisface
\[
 \Ima(p)=U,\qquad\Ker(p)=W,\qquad p^2=p.
\]
Recíprocamente, cualquier endomorfismo idempotente es la proyección
de base $\Ima(p)$ y dirección $\Ker(p)$: se descompone
$v=p(v)+(v-p(v))$ y estos dos subespacios son suplementarios.

Para un endomorfismo $f$, los subespacios de vectores \textbf{fijos}
y \textbf{antifijos} son
\[
 \Fix(f)=\Ker(f-\mathrm{Id}),\qquad
 \Fix^{-}(f)=\Ker(f+\mathrm{Id}).
\]
Una simetría tiene $\Fix(s)=U$, $\Fix^{-}(s)=W$ y $s^2=\mathrm{Id}$.
Recíprocamente, si $s^2=\mathrm{Id}$, entonces
\[
 v=\frac{v+s(v)}2+\frac{v-s(v)}2
\]
descompone $V$ como $\Fix(s)\oplus\Fix^{-}(s)$ y $s$ es la simetría
correspondiente. Sobre $\R$ o $\C$ se puede dividir por $2$.
Las matrices de proyecciones son idempotentes; las de simetrías, involutivas.

Para la misma base y dirección se tiene
\[
 s=2p-\mathrm{Id},\qquad p=\frac{\mathrm{Id}+s}{2}.
\]

\begin{ejemplo}
En $\R^2$, para $U=L((1,0))$ y $W=L((1,1))$,
$(x,y)=(x-y,0)+(y,y)$. Por ello,
$p(x,y)=(x-y,0)$ y $s(x,y)=(x-2y,-y)$.
La dirección $W$ determina qué componente se elimina o se cambia de signo.
\end{ejemplo}

\section{Espacio dual}
\label{sec:dual}

Una \textbf{forma lineal} de $V$ es una aplicación lineal $\varphi:V\to\K$.
El \textbf{espacio dual} es $V^*=\Lin(V,\K)$; si $\dim(V)=n$,
$\dim(V^*)=n$. En una base $B=(v_1,\ldots,v_n)$, una forma se representa
por la matriz \emph{fila}
\[
 \M_B(\varphi)=(\varphi(v_1)\ \cdots\ \varphi(v_n)),\qquad
 \varphi(v)=\M_B(\varphi)\coord{v}{B}.
\]
Una forma no nula es sobreyectiva: si $\varphi(v)=a\ne0$, cualquier
$b\in\K$ es imagen de $(b/a)v$. Su núcleo es un hiperplano.
Recíprocamente, todo hiperplano es núcleo de una forma no nula:
los coeficientes de una ecuación implícita definen esa forma.
La forma nula tiene como núcleo todo $V$.

\subsection*{Base dual y evaluación}

La \textbf{base dual} de $B$ es $B^*=(v_1^*,\ldots,v_n^*)$, definida por
\[
 v_i^*(v_j)=\delta_{ij},\qquad
 \delta_{ij}=\begin{cases}1&i=j,\\0&i\ne j.\end{cases}
\]
Cada $v_i^*$ extrae la coordenada $i$ en $B$. Estas formas son independientes
y, al ser $n$, forman una base del dual. Para cualquier $\varphi\in V^*$,
\[
 \varphi=\sum_{i=1}^n\varphi(v_i)v_i^*,\qquad
 \varphi\Bigl(\sum_{i=1}^nx_iv_i\Bigr)=\sum_{i=1}^n\varphi(v_i)x_i.
\]
Por tanto, las coordenadas de $\varphi$ en $B^*$ son sus valores sobre
los vectores de $B$.

Para calcular una base dual en coordenadas de una base auxiliar $E$,
sea $S=(\coord{v_1}{E}\ \cdots\ \coord{v_n}{E})$.
Las filas que representan las formas $v_i^*$ constituyen la matriz $S^{-1}$:
la condición de evaluación es precisamente $S^{-1}S=I_n$.

\begin{ejemplo}
Para $B=((1,2),(0,1))$ en $\R^2$,
\[
 S=\begin{pmatrix}1&0\\2&1\end{pmatrix},\qquad
 S^{-1}=\begin{pmatrix}1&0\\-2&1\end{pmatrix}.
\]
La base dual es $v_1^*(x,y)=x$, $v_2^*(x,y)=y-2x$.
Así, $\varphi(x,y)=3x+y$ tiene coordenadas $(5,1)$ en $B^*$,
pues $\varphi(1,2)=5$ y $\varphi(0,1)=1$.
\end{ejemplo}
